\documentclass[12pt]{article}
\usepackage[utf8]{inputenc}
\usepackage[T1]{fontenc}
\usepackage{hyperref}
\title{Wayana Language Resources for K-12 Curriculum Development}
\author{Compiled by Claude Sonnet 5 (claude-sonnet-5) for Andr\'as Kornai}
\date{}
\begin{document}
\maketitle
\section{Introduction}
Wayana (ISO 639-3: way; Glottocode waya1269) is a Cariban language of the borderlands of French Guiana, Suriname and Brazil. Ethnologue lists 1,700 users (2012), and about 1,000 speakers were reported along the Maroni in French Guiana. UNESCO lists it as vulnerable.
\section{Primary Educational Resources}
No dictionary, primer or school programme was found. A 2013 article on the education system of the Wayana-Apalai communities of French Guiana was cited in results, but its content was not seen.
\section{Linguistic and Reference Materials}
\begin{itemize}
\item Wikipedia on the language: \url{https://en.wikipedia.org/wiki/Wayana_language}
\item Bakker, Gretenkort and Parkvall (2018), Amerindia 40, on Crevaux's Wayana notes and Wayana Pidgin (historical study): \url{https://amerindia.cnrs.fr/?p=1568}
\item DiCE (Missouri) Wayana PDF: \url{https://dice.missouri.edu/assets/docs/carib/Wayana.pdf}
\item Omniglot: \url{https://mail.omniglot.com/writing/wayana.htm}
\end{itemize}
\section{Religious and Community Resources}
People page: \url{https://en.wikipedia.org/wiki/Wayana}
\section{Digital and Audio Resources}
None found.
\section{Curriculum Development Context}
The language is spread over three countries and little educational material was located. The Aparai lexical site in the results (\url{https://www.webonary.work/apalai/en/?p=9444}) concerns the related Aparai language and is not a Wayana resource.
\end{document}